\subsection{Changement d'état liquide-vapeur de l'eau}

Un système, constitué initialement d'une mole d’eau liquide et une mole d’eau vapeur, maintenu à la température $T = 298 ~ \si{\kelvin}$ et sous la pression $P = P^\circ = 1 ~ \si{\bar}$, est envisagé. 

À cette température, les potentiels chimiques standard de l’eau sont $\mu_{l}^{\star, \circ} = - 237.2 ~ \si{\kilo\joule\per\mole}$ et $\mu_{l}^{\star, \circ} = -228.6 ~ \si{\kilo\joule\per\mole}$.

\Q Montrer que ce système n'est pas initialement en équilibre.

\Q Prévoir son évolution et déterminer l’état final.

\Q La transformation est-elle réversible ?

\Q Calculer l’entropie crée au cours de cette évolution isobare et isotherme.